\documentclass{beamer}

\input{../helper}

\title{Introdução à Computação \\ argumentos de linha de comando e retorno de valor}
\author{Alexandre Rademaker}
\date{}

\begin{document}

\begin{frame}
 \maketitle
\end{frame}

\begin{frame}[fragile,allowframebreaks]{argumentos}
  
  Lembrando que programas CLI podem receber argumentos!

  Podemos carregar o programa e solicitar ao usuário os valores
  desejados.

  Em C, a função mais pode receber dois argumentos: \ilcode{argc}
  (quantidade de argumentos) e \ilcode{argv} (array dos argumentos).

  \framebreak

  Podemos fazer novos programas receber valores da linha de comando!
  Veja o condicional \ilcode{if}!

  \lstinputlisting[style=normalc,caption=argv0.c]{argv0.c}

  \framebreak

  E podemos imprimir todos os argumentos passados

  \lstinputlisting[style=normalc,caption=argv1.c]{argv1.c}

  Repetição com \ilcode{for}!

  \framebreak
  
  Queremos somar um número não pré-determinado de valores.

  Vamos receber da linha de comando a quantidade de números e
  perguntar para o usuário um número por vez.

  \framebreak

  \lstinputlisting[style=tinyc,caption=scores5.c]{scores5.c}

  \framebreak

  Se não passarmos na linha de comando o argumento esperado?

  Se o argumento não for um número que possa ser representado como um
  \ilcode{int}?

\end{frame}

\begin{frame}[fragile,allowframebreaks]{retornando valores}

  Podemos sinalizar que o programa não executou como deveria:

  \lstinputlisting[style=normalc,caption=exit.c]{exit.c}

  \framebreak

  \begin{lstlisting}[style=normalbash]
    > ./exit 12 && echo "Tudo bem"
    hello, 12
    Tudo bem
    > ./exit && echo "Tudo bem"
    Missing arguments
  \end{lstlisting}
    
\end{frame}


\begin{frame}{Exemplos extras}
  Arquivos em \texttt{src/} relacionados ao tema desta semana, para
  estudo complementar (não foram apresentados em aula):
  \begin{itemize}
    \item \texttt{address05.c}
    \item \texttt{argv2.c}
    \item \texttt{arrayspointers.c}
    \item \texttt{parity1.c}
  \end{itemize}
\end{frame}

\end{document}

%%% Local Variables:
%%% mode: latex
%%% TeX-master: t
%%% End:
